\chapter{Retinal Vein Occlusion}
\index{Retinal Vein Occlusion}
\label{sec:Retinal Vein Occlusion}

\section{Overview}
\begin{itemize}[noitemsep]
  \item Retinal vein occlusion (RVO) is a blockage of the veins that drain blood from the retina, causing hemorrhage and fluid leakage
  \item It is the second most common retinal vascular disease after diabetic retinopathy
  \item It typically affects people over 50 and is a major cause of vision loss from macular edema and neovascular complications
\end{itemize}
\section{Causes}
This condition happens because of thrombosis (blood clot) in a retinal vein, most often where the vein crosses a retinal artery (arteriovenous crossing). It is classified as central retinal vein occlusion (CRVO) when the main vein is blocked, or branch retinal vein occlusion (BRVO) when a branch is affected. Contributing factors include hypertension, atherosclerosis, glaucoma, and thrombophilia.
\section{Risk Factors}
\begin{itemize}[noitemsep]
  \item Age over 50 years
  \item Hypertension and high cholesterol
  \item Diabetes
  \item Glaucoma (especially for CRVO)
  \item Smoking
  \item Blood clotting disorders (thrombophilia)
  \item Oral contraceptives and other prothrombotic conditions
\end{itemize}
\section{Signs and Symptoms}
\subsection{Common symptoms}
\begin{itemize}[noitemsep]
  \item Sudden, painless blurring or loss of vision in one eye
  \item Vision loss may be mild and peripheral (BRVO) or severe (CRVO)
  \item Distorted or wavy vision from macular edema
  \item Visual field defects in the affected area
\end{itemize}
\subsection{Emergency warning signs}
\begin{itemize}[noitemsep]
  \item Sudden loss of vision in one eye requires urgent eye examination
  \item New pain, redness, or high eye pressure may indicate neovascular glaucoma
\end{itemize}
\section{Progression and Stages}
Initial symptoms usually develop over hours to days. Two patterns follow: the ischemic (non-perfused) type, which has a poor visual prognosis and high risk of neovascular complications, and the non-ischemic type, which often improves. Macular edema may persist for months to years. Neovascular glaucoma can develop in severe CRVO weeks to months after onset and threatens the eye.
\section{Complications}
\begin{itemize}[noitemsep]
  \item Chronic macular edema with permanent central vision loss
  \item Retinal and iris neovascularization
  \item Neovascular glaucoma (serious, painful, blinding)
  \item Vitreous hemorrhage from new vessels
  \item Retinal detachment in severe cases
  \item Underlying vascular disease progression
\end{itemize}
\section{Diagnosis}
\begin{itemize}[noitemsep]
  \item Dilated fundus examination showing retinal hemorrhages, venous engorgement, and "blood and thunder" appearance in CRVO
  \item Optical coherence tomography (OCT) to assess macular edema
  \item Fluorescein angiography to determine ischemic status
  \item Measurement of intraocular pressure to detect neovascular glaucoma
  \item Blood pressure and cardiovascular risk assessment
  \item Blood tests for clotting disorders in younger patients
\end{itemize}
\section{Prevention}
\begin{itemize}[noitemsep]
  \item Control of blood pressure, cholesterol, and diabetes
  \item Smoking cessation
  \item Regular eye examinations, especially for those with risk factors
  \item Treatment of glaucoma and cardiovascular risk factors
  \item Early evaluation of any sudden monocular visual symptoms
\end{itemize}
\section{Treatment Options}
\begin{itemize}[noitemsep]
  \item Intravitreal anti-VEGF injections (ranibizumab, aflibercept, bevacizumab) for macular edema
  \item Intravitreal corticosteroids (dexamethasone implant) for edema
  \item Grid or scatter laser photocoagulation for BRVO and neovascularization
  \item Panretinal photocoagulation for neovascularization in CRVO
  \item Treatment of neovascular glaucoma with pressure-lowering therapy
  \item Management of systemic cardiovascular risk factors
\end{itemize}
\section{Living with It}
\begin{itemize}[noitemsep]
  \item Attend regular injection and monitoring visits
  \item Control blood pressure, cholesterol, and diabetes
  \item Report new pain, redness, or worsening vision promptly
  \item Use low-vision aids if central vision remains impaired
\end{itemize}
\section{Outlook}
The prognosis is variable. BRVO carries a better outlook, with many eyes recovering spontaneously or with treatment. CRVO has a worse prognosis, especially the ischemic type. With anti-VEGF therapy, many patients with macular edema gain vision. Untreated neovascular glaucoma can cause painful blindness.
